% ==============================================================================
% SECCIÓN 2: METODOLOGÍA Y CICLO DE VIDA
% ==============================================================================
\section{Metodología de Desarrollo y Ciclo de Vida}

\subsection{Modelo Metodológico Ágil}
Para desarrollar la aplicación utilizaremos \textbf{metodología ágil}, integrando los eventos de trabajo de Scrum para la gestión del flujo de trabajo continuo y la monitorización de tareas en curso (\textit{Work In Progress}). Este esquema equilibra la flexibilidad en la asimilación de feedback de usuarios reales con el control formal de alcance, costes, riesgos y plazos de DGP.

\subsection{Ciclo de Vida por Iteraciones}
El proyecto se ejecuta bajo un modelo de \textbf{prototipado evolutivo iterativo}, con 3 iteraciones y una verificación final:

% Revisar fechas
\begin{enumerate}[leftmargin=*]
  \item \textbf{Fase Inicial:} Organizamos el equipo y redactamos este Manual de Coordinación. También preparamos la propuesta técnica con la arquitectura y el presupuesto inicial, definimos los perfiles de usuario y escenarios de uso, y dejamos listos tanto la pila de producto (Product Backlog) como el plan de entregas.
  \item \textbf{Iteración 1:} Configuramos el entorno de trabajo en Flutter y montamos la base de la aplicación. En esta etapa implementamos el inicio de sesión adaptado y la gestión de perfiles accesibles para sacar nuestro primer prototipo funcional y poder probarlo.
  \item \textbf{Iteración 2:} Nos centramos en las actividades del día a día en el centro: los pasos guiados para encargar fotocopias (reprografía), la gestión del comedor con sus menús y dietas especiales, la guía de visitas y un temporizador visual adaptado.
  \item \textbf{Iteración 3:} Añadimos la parte de rutinas personales configurables y el chat accesible para que tutores y alumnos puedan comunicarse. Además, aprovechamos para pulir los colores, contrastes y tamaños de los botones, dejando toda la aplicación estable.
  \item \textbf{Hardening y Verificación Global (Cierre de Iteración 3):} Antes de dar el proyecto por cerrado, hacemos una validación completa con el cliente, ponemos la aplicación al límite para detectar posibles fallos y revisamos que todo cumpla con la protección de datos (RGPD) y la seguridad básica.
  \item \textbf{Defensa Final (12 de enero de 2027):} Exponemos el proyecto y hacemos una demostración en vivo de cómo funciona la aplicación ante los profesores y los representantes del colegio C.E.E. Purísima Concepción.
\end{enumerate}